\subsection{紧空间上的补集与逆极限}

命题 1. --- 设 X 和 Y 是两个拓扑空间，f 是从 X 到 Y 的连续映射。设 \((\mathrm{K}_{\alpha})_{\alpha \in \mathrm{A}}\) 是 X 的紧子集的递减有向族，其交为 K。则 \(f(\mathrm{K}) = \bigcap f(\mathrm{K}_{\alpha})\) 。

事实上，令 \(y\) 是 \(\bigcap_{\alpha \in \mathbf{A}} f(\mathbf{K}_{\alpha})\) 中的一点；对于每个 \(\alpha \in \mathbf{A}\)，集合 \(\mathbf{L}_{\alpha} = \mathbf{K}_{\alpha} \cap \overline{f}^{-1}(y)\) 是紧的且非空。族 \((\mathbf{L}_{\alpha})_{\alpha \in \mathbf{A}}\) 向下有向，因此其交集 \(\mathbf{L}\) 非空。现在，\(\mathbf{L} = \mathbf{K} \cap \overline{f}^{-1}(y)\)，从而 \(y \in f(\mathbf{K})\)。由此我们证明了包含关系 \(f(\mathbf{K}) \supset \bigcap_{\alpha \in \mathbf{A}} f(\mathbf{K}_{\alpha})\)，反向包含显然。

命题 2. --- 给定由 I 索引的拓扑空间逆系统 \((\mathrm{T}_{i}, p_{ij})\)、拓扑空间 T 以及相容且分离的连续映射族 \(p_{i}: T \to T_{i}\)，则：

\begin{enumerate}
\def\labelenumi{\alph{enumi})}
\item
  对于 T 的每个紧子集 K，有 \(\mathbf{K} = \bigcap_{i\in I}\overline{p_i}^{-1}\big(p_i(\mathbf{K})\big)\) 。
\item
  设 K、L 是 T 的两个不交紧子集。存在 \(i \in I\)，使得对于 \(p_{j}(K)\)，\(p_{j}(L)\) 与 \(j \geqslant i\) 不交。
\item
设 \(x\) 是 \(\bigcap_{i \in I} \overline{p_i}^{-1}(p_i(K))\) 中的一点；对于每个 \(i \in I\)，集合 \(K_i\) 由点 \(y\)（属于 \(K\) 且满足 \(p_i(y) = p_i(x)\)）组成，且是 \(K\) 的非空闭子集。若 \(i \leqslant j\)，则 \(K_i \supset K_j\)；由于 \(K\) 紧，集合 \(\bigcap_{i \in I} K_i\) 非空。令 \(y\) 是 \(\bigcap_{i \in I} K_i\) 中的一点；有 \(y \in K\)，且 \(p_i(y) = p_i(x)\) 对所有 \(i \in I\) 成立，从而 \(y = x\)；最后 \(x \in K\)，这证明了包含关系 \(K \supset \bigcap_{i \in I} \overline{p_i}^{-1}(p_i(K))\)；反向包含显然。
\item
  对每个 \(i \in \mathbf{I}\)，置 \(\mathbf{M}_i = \overline{p_i}^{-1}(p_i(\mathbf{K})) \cap \mathbf{L}\)；这是紧空间 \(\mathbf{L}\) 的闭子集，有 \(\mathbf{M}_i \supset \mathbf{M}_j\) 当 \(i \leqslant j\) 时，并且由 \(a\) )，
\end{enumerate}

\[
\bigcap_ {i \in I} M _ {i} = K \cap L = \varnothing .
\]

因此存在指标 \(i\)，使得 \(\mathbf{M}_i = \varnothing\)。对于 \(j \geqslant i\)，有 \(\overline{p_j}^{-1} (p_j(\mathbf{K})) \cap \mathbf{L} = \mathbf{M}_j \subset \mathbf{M}_i = \varnothing\)，从而 \(p_j(\mathbf{K}) \cap p_j(\mathbf{L}) = \varnothing\)。
% semantic-fragments: legacy-input-seam
\relax{}
